\documentclass[tikz,border=6pt]{standalone}
\usepackage{ctex}
\usepackage{amsmath}
\usepackage{amssymb}
\usepackage{pgfplots}
\pgfplotsset{compat=1.18}
\usetikzlibrary{calc,angles,quotes,arrows.meta,positioning,decorations.markings}
\begin{document}
\begin{tikzpicture}

% ===== 左图：单位圆上 sin x >= 1/2 对应的弧段 =====
\begin{scope}[scale=2.0]
  % 坐标轴
  \draw[->,gray!70] (-1.35,0) -- (1.45,0) node[right,black]{$x$};
  \draw[->,gray!70] (0,-1.35) -- (0,1.5) node[above,black]{$y$};
  % 满足条件的弧段填充扇形提示（半透明）
  \fill[green!55!black, opacity=0.18] (0,0) -- (30:1) arc (30:150:1) -- cycle;
  % 单位圆
  \draw[thick] (0,0) circle (1);
  % y = 1/2 水平虚线
  \draw[blue!60!black,dashed] (-1.25,0.5) -- (1.25,0.5);
  \node[blue!60!black,font=\small] at (1.06,0.66) {$y=\tfrac12$};
  % 满足 sin>=1/2 的弧段（加粗绿色）
  \draw[green!50!black, line width=2pt] (30:1) arc (30:150:1);
  % 边界两点
  \coordinate (A) at (30:1);
  \coordinate (B) at (150:1);
  \fill[red!75!black] (A) circle (0.035);
  \fill[red!75!black] (B) circle (0.035);
  \draw[red!75!black,->,thick] (0,0) -- (A);
  \draw[red!75!black,->,thick] (0,0) -- (B);
  % 角度标签
  \node[red!75!black,font=\small] at (1.16,0.98) {$\tfrac{\pi}{6}$};
  \node[red!75!black,font=\small] at (-1.16,0.98) {$\tfrac{5\pi}{6}$};
  % 弧段说明
  \node[green!45!black,font=\small] at (0,1.30) {$\sin x\ge\tfrac12$ 的弧段};
  % 标题
  \node[font=\bfseries] at (0,-1.62) {单位圆：满足条件的上方弧段};
\end{scope}

% ===== 右图：sin 曲线上 y>=1/2 的区间阴影 =====
\begin{scope}[xshift=8.0cm]
  \node[font=\bfseries, anchor=south] at (6.0cm,4.0cm)
    {正弦曲线上 $y\ge\tfrac12$ 的区间};
  \begin{axis}[
      width=12cm, height=6.2cm,
      axis lines=middle,
      xlabel={$x$}, ylabel={$y$},
      xtick={-3.14159,0,3.14159,6.28319,9.42478},
      xticklabels={$-\pi$,$0$,$\pi$,$2\pi$,$3\pi$},
      ytick={-1,0.5,1}, yticklabels={$-1$,$\tfrac12$,$1$},
      xmin=-3.8, xmax=10.4, ymin=-1.55, ymax=1.5,
      enlargelimits=false,
      tick label style={font=\small},
      label style={font=\small},
      legend style={at={(0.5,-0.14)}, anchor=north, legend columns=2,
                    font=\small, draw=none, fill=none, column sep=1.2em},
    ]
    % 阴影：在每个解区间内填充 sin 曲线与 y=1/2 之间的区域
    \foreach \off in {-6.283185, 0, 6.283185} {
      \pgfmathsetmacro{\xl}{0.523599+\off}  % pi/6
      \pgfmathsetmacro{\xr}{2.617994+\off}  % 5pi/6
      \edef\temp{\noexpand\addplot[draw=none, fill=green!55!black, fill opacity=0.22,
        domain=\xl:\xr, samples=60] {sin(deg(x))} \noexpand\closedcycle;}
      \temp
    }
    % sin 曲线
    \addplot[domain=-3.8:10.4, samples=300, thick, black] {sin(deg(x))};
    \addlegendentry{$y=\sin x$}
    % y = 1/2 水平线
    \addplot[domain=-3.8:10.4, samples=2, blue!60!black, dashed, thick] {0.5};
    \addlegendentry{$y=\tfrac12$}
    % 区间端点标记
    \foreach \off in {-6.283185, 0, 6.283185} {
      \pgfmathsetmacro{\xl}{0.523599+\off}
      \pgfmathsetmacro{\xr}{2.617994+\off}
      \edef\temp{\noexpand\addplot[only marks, mark=*, red!75!black, mark size=1.5pt]
        coordinates {(\xl,0.5) (\xr,0.5)};}
      \temp
    }
  \end{axis}

  % ===== 数轴：解集区间 =====
  \begin{scope}[yshift=-3.55cm]
    \draw[->,gray!70] (0,0) -- (11.2,0) node[right,black]{$x$};
    % 主刻度
    \foreach \xc/\lab in {1.0/$-\pi$, 3.1/$0$, 5.2/$\pi$, 7.3/$2\pi$, 9.4/$3\pi$} {
      \draw[gray!70] (\xc,0.08) -- (\xc,-0.08);
      \node[font=\small,below] at (\xc,-0.10) {\lab};
    }
    % 解集区间（与上图三段对应，闭区间）
    % 像素映射：pi -> 2.1 单位/π，0 在 x=3.1
    % 区间 [pi/6,5pi/6]: 中心区段
    \foreach \kc in {-1,0,1} {
      \pgfmathsetmacro{\xa}{3.1 + 2.1*(0.16667 + 2*\kc)} % pi/6
      \pgfmathsetmacro{\xb}{3.1 + 2.1*(0.83333 + 2*\kc)} % 5pi/6
      \draw[green!45!black, line width=2.6pt] (\xa,0.0) -- (\xb,0.0);
      \fill[green!45!black] (\xa,0) circle (0.06);
      \fill[green!45!black] (\xb,0) circle (0.06);
    }
    \node[green!40!black,font=\small,fill=white,fill opacity=0.8,text opacity=1,
          inner sep=2pt, anchor=north] at (5.5,-0.55)
      {解集：$\displaystyle\Big[\tfrac\pi6+2k\pi,\ \tfrac{5\pi}6+2k\pi\Big],\quad k\in\mathbb{Z}$};
  \end{scope}
\end{scope}

\end{tikzpicture}
\end{document}
